\documentclass[11pt]{article}
\usepackage[T1]{fontenc}
\usepackage[margin=1in]{geometry}
\usepackage{amsmath,amssymb,amsthm}
\usepackage[colorlinks=true,linkcolor=blue,citecolor=blue,urlcolor=blue]{hyperref}
\newtheorem{theorem}{Theorem}
\newtheorem{lemma}[theorem]{Lemma}
\newcommand{\R}{\mathbb R}
\newcommand{\pdisc}{\operatorname{pdisc}}
\title{Kintali's counterexample to fixed-order prefix balancing}
\author{Kaiwen Jin\\\small\href{mailto:yui.ui.siki@gmail.com}{yui.ui.siki@gmail.com}}
\date{October 12, 2026}

\begin{document}
\maketitle
\begin{abstract}
The strong Koml\'os question asks whether one signing can keep every
prefix of any fixed ordered family of Euclidean-unit vectors bounded
in the maximum norm by an absolute constant. We give a self-contained
exposition of Kintali's counterexample. His finite recursive family has
dimension $d_k=2^{2^k-1}$ and prefix discrepancy at least
$2^{k/2-1}=\tfrac12\sqrt{1+\log_2d_k}$ for every signing.
The proof freezes a large mean in each of several coordinate pairs and
then forces a large difference in one of those pairs. All witness choices
are represented in a fixed ordered word, so the argument applies to
signs chosen with knowledge of the entire input. The construction and
amplification argument are due to Kintali.
\end{abstract}

\section{The question and its source}

For a matrix $A=(v_1,\ldots,v_T)$ with ordered columns in $\R^d$, define
\[
\pdisc(A)=\min_{\varepsilon\in\{-1,1\}^T}
\max_{1\le t\le T}\left\|\sum_{i=1}^t\varepsilon_i v_i\right\|_\infty,
\qquad \|x\|_\infty=\max_{1\le r\le d}|x_r|.
\]
The order is fixed, and the same signs must work for all prefixes.
They may depend on the entire input. The strong Koml\'os conjecture asks
whether $\pdisc(A)$ has an absolute upper bound whenever
$\|v_i\|_2\le1$. Bansal et al.~\cite[Open Problem 6.1]{bansal} state
this question, which is the target of AIM problem 325~\cite{aim}.
It differs from the terminal-sum problem: a small value of
$\|\sum_i\varepsilon_i v_i\|_\infty$ allows large intermediate
excursions. Karingula and Lovett's terminal-sum result and their
Conjecture 1.6 explicitly retain this distinction~\cite{karingula}.

Kintali's preprint~\cite[Lemma 2.2 and Theorem 3.1]{kintali} gives
the following negative answer. This note explains its construction and
proof; it makes no claim of a new construction or mathematical priority.

\begin{theorem}[Kintali]\label{thm:main}
For every integer $k\ge0$, there is a finite fixed ordered matrix
$A_k\in\R^{d_k\times n_k}$, where $d_k=2^{2^k-1}$ and $n_k\ge1$,
whose columns have Euclidean norm one, such that every
$\varepsilon\in\{-1,1\}^{n_k}$ satisfies
\begin{equation}\label{eq:main}
\max_{1\le t\le n_k}\left\|\sum_{i=1}^t
\varepsilon_i(A_k)_i\right\|_\infty
\ge 2^{k/2-1}=\frac12\sqrt{1+\log_2d_k}.
\end{equation}
Consequently there is no absolute constant in the strong Koml\'os question.
\end{theorem}

Here $(A_k)_i$ is column $i$, and $n_k$ is its finite number of columns.
The right side grows without bound. The length of the construction is
large; neither an efficient algorithm nor an optimal lower bound is asserted.

\section{The construction and amplification}

Copies of a matrix can begin after earlier copies have moved the current
sum. We therefore use a statement that allows an arbitrary initial offset.
Put $E=\{c\in\R:|c|\ge1\}$. For an ordered $d\times n$ matrix
$M=(M_1,\ldots,M_n)$ and $a\ge0$, write $R(M,a)$ for
\begin{equation}\label{eq:robust}
\forall c\in E^n\ \forall u\in\R^d:
\quad\max_{0\le t\le n}\left\|u+\sum_{j=1}^t c_jM_j\right\|_\infty\ge a.
\end{equation}
The prefix with $t=0$ is included in this auxiliary property.

Assume $d,n\ge1$. In dimension $2d^2$, index coordinates by
$(\ell,r,\sigma)$ for $\ell,r\in[d]=\{1,\ldots,d\}$ and
$\sigma\in\{+,-\}$, and let $e_{\ell,r,\sigma}$ be the corresponding
standard vector. For each source column $j$, define
\begin{align}
O_{\ell,j}&=\sum_{r=1}^d m_{rj}
 (e_{\ell,r,+}+e_{\ell,r,-}),\label{eq:original}\\
D_{\rho,j}&=\sum_{\ell=1}^d m_{\ell j}
 (e_{\ell,r_\ell,+}-e_{\ell,r_\ell,-}),
 \qquad \rho=(r_1,\ldots,r_d)\in[d]^d.\label{eq:detector}
\end{align}
The source row in the second formula is $\ell$; $r_\ell$ chooses only
the physical coordinate pair. Different depths $\ell$ use disjoint pairs.
An original column will change the mean of a pair, while a detector
column will change its difference.

Let $Q$ concatenate $(D_{\rho,1},\ldots,D_{\rho,n})$ for every
$\rho\in[d]^d$ in lexicographic order. Define words backwards by
\begin{equation}\label{eq:word}
W_{d+1}=Q,\qquad
W_\ell=(W_{\ell+1},O_{\ell,1},W_{\ell+1},\ldots,
 O_{\ell,n},W_{\ell+1})\quad(\ell=d,\ldots,1).
\end{equation}
The matrix $F(M)$ has ordered column word $W_1$. Every physical occurrence
has its own coefficient, even when two occurrences are equal as vectors.
The entire word is specified before these coefficients are chosen.

Each output column corresponding to $M_j$ has squared norm
$2\|M_j\|_2^2$. This follows by counting its two copies of each source
entry in distinct coordinates in~\eqref{eq:original} and
\eqref{eq:detector}. In particular, $s$ nonzero entries in $\{-1,1\}$
per source column become exactly $2s$ such entries per output column.
The output dimension and length are
\begin{equation}\label{eq:size}
d'=2d^2,\qquad n'=(1+nd^d)(n+1)^d-1.
\end{equation}
Indeed $|Q|=nd^d$ and $|W_\ell|+1=(n+1)(|W_{\ell+1}|+1)$.
Thus the construction is finite for every finite source matrix.

\begin{lemma}[Kintali]\label{lem:amplify}
If $R(M,a)$ holds, then $R(F(M),2a)$ holds.
\end{lemma}

\begin{proof}
Fix arbitrary coefficients in $E$ on all $n'$ physical occurrences and
an arbitrary initial vector $w\in\R^{2d^2}$. Let $p(t)$ be the resulting
path, including $p(0)=w$. For each coordinate pair set
\[
h_{\ell,r}(t)=\frac{p_{\ell,r,+}(t)+p_{\ell,r,-}(t)}2,
\qquad z_{\ell,r}(t)=\frac{p_{\ell,r,+}(t)-p_{\ell,r,-}(t)}2.
\]
Thus $p_{\ell,r,\pm}=h_{\ell,r}\pm z_{\ell,r}$. An occurrence of
$O_{\ell,j}$ changes only the depth-$\ell$ mean vector, by its
coefficient times $M_j$. No original changes any $z$. A detector
$D_{\rho,j}$ changes $z_{\ell,r_\ell}$ by its coefficient times
$m_{\ell j}$ for every $\ell$, and changes no mean.

In an occurrence of $W_\ell$, let $h^0$ be the depth-$\ell$ mean
vector at its start and $\alpha_j\in E$ the coefficients of its $n$
separating original columns. At the starts of its $n+1$ child occurrences
of $W_{\ell+1}$ the mean vector has values
\[
h^0+\sum_{j=1}^s\alpha_jM_j,\qquad 0\le s\le n.
\]
Property~\eqref{eq:robust} selects a child and a row $r_\ell$ with
$|h_{\ell,r_\ell}|\ge a$. That mean stays fixed throughout the
chosen child, since it contains only deeper originals and detectors.
Starting at $W_1$, repeat this selection inside each chosen child.
After $d$ selections there is one occurrence of $Q$ and rows
$r_1,\ldots,r_d$ such that all means $h_{\ell,r_\ell}$ are constant
there and have absolute value at least $a$. Previously selected means
remain fixed on every descendant interval.

This occurrence of $Q$ contains the detector word for
$\rho=(r_1,\ldots,r_d)$. If its occurrence coefficients are
$\beta_1,\ldots,\beta_n$, the vector of selected differences follows
\[
(z_{1,r_1},\ldots,z_{d,r_d})_{\rm start}
 +\sum_{j=1}^s\beta_jM_j,\qquad 0\le s\le n.
\]
Its initial value may reflect earlier detectors. Property
\eqref{eq:robust} permits precisely this arbitrary offset, and gives a
time and depth $\ell$ with $|z_{\ell,r_\ell}|\ge a$. At the same
time $|h_{\ell,r_\ell}|\ge a$. For real $h,z$,
\[
\max\{|h+z|,|h-z|\}=|h|+|z|,
\]
as follows by squaring both sides. One physical coordinate therefore has
magnitude at least $2a$ at a global prefix time.
The initial vector and all occurrence coefficients were arbitrary.
Choosing the witness interval and tuple after fixing them changes
neither the input word nor its coefficients. This proves the lemma.
\end{proof}

\section{Iteration and unit normalization}

\begin{proof}[Proof of Theorem~\ref{thm:main}]
Set $M_0=[1]$ and $M_{k+1}=F(M_k)$. For $u\in\R$ and $|c|\ge1$,
\[
|c|\le|u|+|u+c|,
\]
so $R(M_0,1/2)$ holds. Lemma~\ref{lem:amplify} gives
$R(M_k,2^{k-1})$ for every integer $k\ge0$.
Each $M_k$ is an integer matrix with exactly $s_k$ nonzero entries in
$\{-1,1\}$ per column. The norm count and~\eqref{eq:size} give
\[
s_0=d_0=n_0=1,\qquad s_{k+1}=2s_k,\qquad d_{k+1}=2d_k^2,
\]
and the finite positive length obeys
\[
n_{k+1}=(1+n_kd_k^{d_k})(n_k+1)^{d_k}-1.
\]
Induction yields $s_k=2^k$ and $d_k=2^{2^k-1}$.
Define $A_k=M_k/\sqrt{s_k}$. Every column then has squared Euclidean
norm $s_k/s_k=1$. For any signing, apply $R(M_k,2^{k-1})$ with
offset zero and divide by $\sqrt{s_k}$. The result is~\eqref{eq:main},
since $2^{k-1}/\sqrt{2^k}=2^{k/2-1}$ and
$1+\log_2d_k=2^k$. The zero prefix has norm zero, so it cannot be
the witness to this positive bound.

For any proposed $C>0$, choose a positive integer $m>2C$ and take
$k=2m$. Since $2^m\ge m+1>2C$, the bound $2^{k/2-1}$ exceeds $C$.
The fixed finite family $A_k$ then defeats every signing. Every column is
retained, the order is unchanged, and the signs may use the entire family.
\end{proof}

\paragraph{AI-use statement.}
OpenAI Codex assisted with source retrieval, reconstruction and checking
of the argument, and preparation of this exposition. The construction
and proof mechanism are credited to Shiva Kintali. The checking performed
for this note was AI-assisted; no independent human review is claimed.

\begin{thebibliography}{4}
\bibitem{bansal}
N. Bansal, H. Jiang, R. Meka, S. Singla, and M. Sinha,
\emph{Prefix Discrepancy, Smoothed Analysis, and Combinatorial Vector Balancing},
ITCS 2022, LIPIcs 215, 13:1--13:22.
\href{https://doi.org/10.4230/LIPIcs.ITCS.2022.13}{doi:10.4230/LIPIcs.ITCS.2022.13}.
Full version: \href{https://arxiv.org/abs/2111.07049v1}{arXiv:2111.07049v1},
Section 6, Open Problem 6.1.

\bibitem{karingula}
S. R. Karingula and S. Lovett,
\emph{An elementary proof of the Koml\'os conjecture},
\href{https://eccc.weizmann.ac.il/report/2026/188/}{ECCC TR26-188}, 2026.
\href{https://arxiv.org/abs/2609.20979v1}{arXiv:2609.20979v1},
Theorem 1.2 and Conjecture 1.6.

\bibitem{kintali}
S. Kintali, \emph{On the Fixed-Order Strong Koml\'os Conjecture},
\href{https://eccc.weizmann.ac.il/report/2026/242/}{ECCC TR26-242}, 2026,
Sections 2--3.1, Lemma 2.2 and Theorem 3.1, pp. 2--5.
The PDF title page is dated September 15; the ECCC report record is dated
September 26; its publication field gives October 9.

\bibitem{aim}
\emph{AIM---Open Applied Problems}, problem 325,
\href{https://github.com/MColbrook/AIM/blob/c9929805cece044705abaaebe1cebddac1ba11a6/problems/325-strong-komlos-prefix-discrepancy.md}{fixed-order prefix discrepancy statement},
revision \texttt{c9929805cece044705abaaebe1cebddac1ba11a6},
accessed October 12, 2026.
\end{thebibliography}
\end{document}
